\scanpage{9}

signifie ceci : \deleted{Corollaire} \uline{Toute déformation \unsure{infinitésime|infinitésimale} du premier ordre de $C$ dans $A$} [si $C$ est plongé \unsure{jacobiennement|provisoirement|positionnée} ds $A$] \uline{est induit par une translation infinitésimale de $A$} (qui est bien déterminée d'ailleurs, puisque $C$ n'a pas d'automorphismes \inserted{infinitésimaux}).

Or \deleted{comme} les déformations infinitésimales en question sont définies par un élément de $H^0(C,N)$, où $N$ est le faisceau normal de $C$ dans $A$. Considérons la suite exacte
\diagram{\begin{tikzcd}[ampersand replacement=\&,column sep=small]
0\arrow[r,dash]\&\underline T\arrow[r]\&\underline V\arrow[r,dash]\&\underline N\arrow[r]\&0
\end{tikzcd}}
où $V$ est induit par le faisceau tangent à $A$ donc constant de rang $g$, $T$ est le faisceau tangent à $C$ donc d$^\circ$ inversible de degré $2(1-g)<0$, \deleted{et} \unsure{où|on} a donc $H^0(C,\underline T)=0$, $H^0(\unsure{C|O},\underline V)\simeq k^g$ $\underset{\unsure{\sim}}{=}$ \deleted{$\unsure{H^0(C}$} $\mathrm{Im}\,H^0(A,\unsure{\mathscr T|\mathscr J|\check{\Omega}}_{A/k})$, et la question est donc \unsure{exactement|seulement} si $H^0(C,\underline V)\mathbin{\unsure{-|\to}}H^0(\underline C,\underline N)$ est \uline{surjectif} i.e. $H^1(C,\underline T)\to H^1(C,\underline V)\xleftarrow{\sim}\deleted{\unsure{k}}H^1(A,\mathscr T_{A/k})$ injectif, i.e.

% archive sheet number: 8
